\begin{table}[htbp]
\centering
\footnotesize
\begin{tabularx}{\textwidth}{L{3.2cm}YYYY}
\toprule
\textbf{Run} & \textbf{1 Physical (max)} & \textbf{2a Spearman $\geq$ 0.49} & \textbf{2b Gaps $\leq$ 2.0 / 1.0 / 5.75} & \textbf{4 Level, CTD within} \\
\midrule
Conditional, test period & yes (3) & yes (0.5099) & no (1.1 / 1.7 / 6.8) & no (8.0) \\
Conditional, training period & yes (3) & no (0.4898) & yes (0.8 / 1.0 / 5.2) & periods (3.9) \\
Generative, test period & yes (3) & no (0.4483) & no (1.7 / 1.5 / 4.0) & no (11.2) \\
\midrule
\multicolumn{5}{l}{\textit{3 Extreme condition, TTT median / mean at arrival intensity 0.8, 1.0, 1.2}} \\
Conditional, test period & \multicolumn{4}{l}{26.0 / 30.4 $\rightarrow$ 27.8 / 33.1 $\rightarrow$ 33.8 / 44.5 min, rises} \\
\bottomrule
\end{tabularx}
\par\smallskip{\footnotesize\raggedright Criteria of Section~\ref{sec:exp_e4_runs}. The thresholds come from the differences between real periods. The Spearman threshold is that of the weakest test week. The thresholds of the profile gaps (approach, stop duration and TTT by trucks in the terminal, in minutes) are the largest differences between weeks. The CTD must lie within the differences between weeks (3.1 minutes) or between periods (5.5 minutes).\par}
\caption{Acceptance criteria for the final configuration.}
\label{tab:rq2_acceptance}
\end{table}
